\chapter{C2: Fortgeschrittene Grammatik}
\label{cha:C2Fortgeschritten}

Dieses Kapitel behandelt fortgeschrittene grammatikalische Themen für das C2-Niveau.

\section{Komplexe Satzkonstruktionen}
\label{sec:C2KomplexeSaetze}

\subsection{Parenthesen (Einschübe)}
Parenthesen sind erklärende Einschübe im Satz:

\begin{itemize}
    \item \textbf{Relativsatz als Parenthese}: Der Mann – ich kenne ihn seit Jahren – arbeitet bei uns.
    \item \textbf{Nomenphrase als Parenthese}: Das Auto, um das es geht, ist quite teuer.
    \item \textbf{Adverbialphrase als Parenthese}: Ich werde, ehrlich gesagt, daran zweifeln.
\end{itemize}

\subsection{Attributive Nebensätze}
\begin{itemize}
    \item \textbf{Nebensatz}: Das ist das Buch, das ich gekauft habe.
    \item \textbf{Attribut}: Das gekaufte Buch (Partizipialattribut)
    \item \textbf{Erweitertes Partizip}: das gestern in der Bibliothek ausgeliehene Buch
\end{itemize}

\section{Infinitivgruppen}
\label{sec:C2Infinitiv}

\begin{center}
\begin{tabular}{|l|l|p{6cm}|}
\hline
Typ & Beispiel & Bedeutung \\ \hline
um zu & Ich lerne, um eine Prüfung zu bestehen. & Zweck \\ \hline
ohne zu & Er ging, ohne ein Wort zu sagen. & Negativer Begleitumstand \\ \hline
statt zu & Sie blieb zu Hause, statt zu arbeiten. & Alternatives \\ \hline
ums zu & Es geht ums Geld zu sparen. & Zweck (gehoben) \\ \hline
\end{tabular}
\end{center}

\section{Modales Passiv}
\label{sec:C2ModalPassiv}

\begin{center}
\begin{tabular}{|l|c|c|}
\hline
Modalität & Aktiv & Passiv \\ \hline
müssen & Er muss das machen. & Das muss gemacht werden. \\ \hline
können & Man kann das machen. & Das kann gemacht werden. \\ \hline
sollen & Die Arbeit soll gemacht werden. & - \\ \hline
wollen & Er will das gemacht wissen. & Das soll gemacht werden. (wollen-Passiv) \\ \hline
dürfen & Das darf nicht gemacht werden. & - \\ \hline
\end{tabular}
\end{center}

\section{Konjunktiv I und II in der gehobenen Sprache}
\label{sec:C2KonjunktivGehoben}

\bigskip
\noindent\textbf{Konjunktiv I als Indirekte Rede:}

\begin{itemize}
    \item \textbf{Nach Verben des Sagens/Denkens}: Er sagte, er sei krank. (Original: „Ich bin krank.")
    \item \textbf{Nach „man hört/man munkelt"}: Man sagt, er sei ein Genie.
    \item \textbf{Zeitungsstil}: Die Ministerin teilte mit, die Verhandlungen seien erfolgreich abgeschlossen.
\end{itemize}

\bigskip
\noindent\textbf{Konjunktiv II für Höflichkeit und Irrealität:}

\begin{center}
\begin{tabular}{|l|c|c|}
\hline
Form & Verwendung & Beispiel \\ \hline
wäre + Partizip II & Irreale Bedingung & Wäre ich früher gekommen, hätte ich ihn getroffen. \\ \hline
hätte + Partizip II & Irreale Vergangenheit & Hättest du angerufen, wäre ich zu Hause geblieben. \\ \hline
würde + Infinitiv & Höfliche Bitte & Würden Sie mir bitte helfen? \\ \hline
\end{tabular}
\end{center}

\chapterend

\chapter{C2: Stilistische Feinheiten}
\label{cha:C2Stil}

\section{Nominalstil vs. Verbalstil}
\label{sec:C2Nominalverbal}

\bigskip
\begin{center}
\begin{tabular}{|l|l|l|}
\hline
Begriff & Nominalstil (formell) & Verbalstil (informell) \\ \hline
Entscheidung treffen & eine Entscheidung treffen & entscheiden \\ \hline
in Betracht ziehen & in Betracht ziehen & erwägen \\ \hline
zur Kenntnis nehmen & zur Kenntnis nehmen & wissen \\ \hline
in Frage stellen & in Frage stellen & bezweifeln \\ \hline
durchführen & zur Durchführung bringen & durchführen \\ \hline
\end{tabular}
\end{center}

\section{Partizipien als Attribute}
\label{sec:C2PartizipAttribute}

\bigskip
\begin{center}
\begin{tabular}{|l|c|c|}
\hline
Partizip & aktiv & passiv \\ \hline
Präsens & die lesende Frau & - \\ \hline
Perfekt & - & das gelesene Buch \\ \hline
Futur & die zu lesende Aufgabe & - \\ \hline
\end{tabular}
\end{center}

\bigskip
\textbf{Erweiterte Attribute}:
\begin{itemize}
    \item das in der Bibliothek ausgeliehene Buch
    \item der von allen respektierte Lehrer
    \item die noch zu erledigenden Aufgaben
\end{itemize}

\section{Wortstellungsvarianten}
\label{sec:C2Wortstellung}

\begin{itemize}
    \item \textbf{Nachfeld}: Ich habe heute, wie jeden Tag, keine Zeit.
    \item \textbf{Vorfeld}: Morgen werde ich, wenn das Wetter gut ist, spazieren gehen.
    \item \textbf{Mittelfeld}: Ich habe ihm gestern das Buch gegeben.
    \item \textbf{Herausstellung}: Was das Essen betrifft, bin ich sehr kritisch.
\end{itemize}

\section{Konjunktiv II: Feine Unterschiede}
\label{sec:C2KonjunktivFein}

\bigskip
\begin{center}
\begin{tabular}{|l|l|p{6cm}|}
\hline
Form & Verwendung & Beispiel \\ \hline
wäre/wär & \begin{minipage}{6cm}Vermutung, spekulative Annahme\end{minipage} & Er wäre krank (aber ich bin nicht sicher) \\ \hline
hätte & \begin{minipage}{6cm}Irreale Vergangenheit, Bedauern\end{minipage} & Hätte ich das gewusst! \\ \hline
würde & \begin{minipage}{6cm}Höfliche Bitte, indirekte Aussage\end{minipage} & Würden Sie mir bitte helfen? \\ \hline
täte & \begin{minipage}{6cm}Ersetzt „würde + Verb" (gehoben)\end{minipage} & Er täte gut daran zu kommen. \\ \hline
\end{tabular}
\end{center}

\section{Subjunktoren auf Niveau C2}
\label{sec:C2Subjunktoren}

\bigskip
\begin{center}
\begin{tabular}{|l|l|p{6cm}|}
\hline
Subjunktor & Bedeutung & Beispiel \\ \hline
insofern ... als & Einschränkung & Er ist klug, insofern als er viel liest. \\ \hline
sofern & Bedingung & Ich komme, sofern du Zeit hast. \\ \hline
je ... desto & Proportionalität & Je mehr, desto besser. \\ \hline
als ob/als wenn & Vergleich & Er tut, als ob er everything wüsste. \\ \hline
ohne dass & Begleitumstand (negativ) & Er ging, ohne dass jemand es bemerkte. \\ \hline
außer dass & Ausnahme & Ich mag ihn, außer dass er zu laut ist. \\ \hline
geschweige denn & Steigerung (negativ) & Er kann nicht lesen, geschweige denn schreiben. \\ \hline
\end{tabular}
\end{center}

\section{Feste Wendungen und Redewendungen}
\label{sec:C2Wendungen}

\bigskip
\noindent\textbf{Idiomatische Ausdrücke auf C2-Niveau:}

\begin{itemize}
    \item \textbf{Nominalisiert}: anhand, aufgrund, infolge, zufolge, zugunsten/zuungunsten
    \item \textbf{Verbale Wendungen}: es handelt sich um, zum Ausdruck bringen, in Betracht ziehen, zur Anwendung kommen
    \item \textbf{Partikelverben}: durch die Finger gehen, sich durchsetzen, auf den Grund gehen, außer Acht lassen
\end{itemize}

\bigskip
\noindent\textbf{Modalpartikeln auf C2-Niveau:}

\begin{itemize}
    \item \textbf{ja}: Ich kenne ihn ja.
    \item \textbf{doch}: Das ist doch klar!
    \item \textbf{wohl}: Das wird wohl stimmen.
    \item \textbf{bloß/nur}: Ich will bloß/nur wissen...
    \item \textbf{immerhin}: Es ist immerhin besser als nichts.
    \item \textbf{gar}: Das ist gar nicht so schwer.
\end{itemize}

\section{Textsorten und Register}
\label{sec:C2Register}

\bigskip
\begin{center}
\begin{tabular}{|l|p{5cm}|p{5cm}|}
\hline
Stil & Merkmale & Beispiel \\ \hline
Wissenschaftlich & Nominalstil, Passiv, unpersönlich & „Es wird davon ausgegangen..." \\ \hline
Journalistisch & Kurz, prägnant, aktivierend & „Die Regierung kündigte an..." \\ \hline
Literarisch & Bildhaft, emotional, subjektiv & „Der Himmel hing low über der Stadt..." \\ \hline
Bürokratisch & Formell, komplex, standardisiert & „Hiermit wird bestätigt..." \\ \hline
Gesprochen & Umgangssprachlich, elliptisch & „Isch hab kenne Zeit." \\ \hline
\end{tabular}
\end{center}

\chapterend

\chapter{C2: Spezialthemen}
\label{cha:C2Spezial}

\section{Verben mit Präpositionen (C2)}
\label{sec:C2VerbPraep}

\bigskip
\begin{center}
\begin{tabular}{|l|l|l|}
\hline
Verb & Präposition & Beispiel \\ \hline
absehen von & Akk & Ich sehe von einer Antwort ab. \\ \hline
achten auf & Akk & Achten Sie auf die Qualität! \\ \hline
ankommen auf & Akk & Es kommt auf das Ergebnis an. \\ \hline
sich ängstigen vor & Dat & Ich ängstige mich vor dem Resultat. \\ \hline
sich aufhalten in & Dat & Er hält sich in Berlin auf. \\ \hline
sich bedanken bei/für & Dat/Akk & Ich bedanke mich bei Ihnen für die Hilfe. \\ \hline
beginnen mit & Dat & Die Sitzung beginnt mit einer Rede. \\ \hline
berichten über & Akk & Er berichtet über das Ereignis. \\ \hline
bestehen aus & Dat & Das Team besteht aus fünf Mitgliedern. \\ \hline
bestehen auf & Dat & Er besteht auf seiner Meinung. \\ \hline
sich bewerben um & Akk & Sie bewirbt sich um die Stelle. \\ \hline
sich einlassen auf & Akk & Er lässt sich auf das Experiment ein. \\ \hline
sich erkundigen nach & Dat & Ich erkundige mich nach dem Weg. \\ \hline
sich freuen auf/über & Akk & Ich freue mich auf die Reise. \\ \hline
gehören zu & Dat & Das gehört zur Aufgabe. \\ \hline
sich halten an & Akk & Bitte halten Sie sich an die Regeln! \\ \hline
hindern an & Dat & Ich wurde daran gehindert zu kommen. \\ \hline
hoffen auf & Akk & Ich hoffe auf eine positive Antwort. \\ \hline
sich informieren über & Akk & Informieren Sie sich über die Konditionen. \\ \hline
klagen über & Akk & Er klagt über Kopfschmerzen. \\ \hline
sich konzentrieren auf & Akk & Ich konzentriere mich auf die Aufgabe. \\ \hline
leiden an & Dat & Er leidet an einer Krankheit. \\ \hline
rechnen mit & Dat & Ich rechne mit deiner Unterstützung. \\ \hline
sich richten nach & Dat & Das richtet sich nach den Umständen. \\ \hline
schreiben an & Akk & Er schreibt an einen Brief. \\ \hline
sprechen über/von & Akk & Wir sprechen über das Thema. \\ \hline
stimmen für/gegen & Akk & Die Mehrheit stimmt für den Vorschlag. \\ \hline
sich streiten um & Akk & Sie streiten um das Erbe. \\ \hline
sich kümmern um & Akk & Ich kümmere mich um das Problem. \\ \hline
trauen auf & Akk & Er traut auf sein Glück. \\ \hline
vertrauen auf & Akk & Ich vertraue auf deine Hilfe. \\ \hline
warnen vor & Dat & Die Polizei warnt vor dem Betrüger. \\ \hline
wetten auf & Akk & Er wettet auf den Sieg. \\ \hline
wissen um & Akk & Er weiß um die Bedeutung. \\ \hline
zweifeln an & Dat & Ich zweifele an seiner Kompetenz. \\ \hline
\end{tabular}
\end{center}

\section{Weitere Präpositionen mit Dativ/Akkusativ}
\label{sec:C2WeiterePraep}

\bigskip
\begin{center}
\begin{tabular}{|l|l|p{6cm}|}
\hline
Präposition & Fall & Verwendung \\ \hline
ab (zeitlich) & Dat & Ab morgen ist er nicht mehr hier. \\ \hline
ab (räumlich) & Dat & Ab dem Bahnhof sind wir zu Fuß. \\ \hline
außer & Dat & Außer mir war niemand da. \\ \hline
aus (Ursache) & Dat & Aus Neid hat er es getan. \\ \hline
bei (Gelegenheit) & Dat & Bei dieser Gelegenheit möchte ich... \\ \hline
bis & Akk & Bis nächsten Montag. \\ \hline
durch (zeitlich) & Akk & Den ganzen Tag durch. \\ \hline
gegenüber & Dat & Gegenüber dem Rathaus. \\ \hline
mit & Dat & Mit 50 Jahren. \\ \hline
nach & Dat & Nach dem Essen. \\ \hline
seit & Dat & Seit einer Woche. \\ \hline
von & Dat & Von Montag bis Freitag. \\ \hline
zu & Dat & Zur Zeit der Abreise. \\ \hline
\end{tabular}
\end{center}

\section{Nominalisierung im C2-Kontext}
\label{sec:C2Nominalisierung}

\bigskip
\noindent\textbf{Häufige Nominalisierungen:}

\begin{center}
\begin{tabular}{|l|l|}
\hline
Verb/Nomen & Nominalisierung \\ \hline
entscheiden & die Entscheidung treffen \\ \hline
erklären & eine Erklärung abgeben \\ \hline
diskutieren & zur Diskussion stellen \\ \hline
beschließen & einen Beschluss fassen \\ \hline
vereinbaren & eine Vereinbarung treffen \\ \hline
versprechen & ein Versprechen geben \\ \hline
ablehnen & die Ablehnung aussprechen \\ \hline
zustimmen & die Zustimmung erteilen \\ \hline
\end{tabular}
\end{center}

\chapterend